\section{Sequence Synchronization Procedure}
\label{app:sequence-synchronization}

This appendix describes how raw per-channel CSV data become the synchronized
multi-channel tensors $\tensor{x}_{u\,v}(t)$ and $\tensor{f}(t)$ for an edge
$(u,v)$.  The important point is
that synchronization happens in two stages: first each channel is regularized
onto its own key lattice, and only then are channels assembled together around
a common anchor key.

\subsection{Per-Channel Lattice Regularization}
\label{appsub:per-channel-lattice-regularization}

Each source channel is processed independently before batching.  The CSV reader
first infers a regular key increment from the source file, then scans the file
in nondecreasing key order.  During that pass:
\begin{enumerate}
\item Duplicate or near-duplicate keys are collapsed by keeping the latest
record.
\item If the gap to the next valid record spans several regular increments, the
current anchor record is emitted and the missing intermediate key values are
filled with null records.
\item For kline data, malformed key values may be rewritten through the
channel-specific lattice-recovery logic if they can be brought back onto the
regular grid; otherwise the file is rejected.
\end{enumerate}

The result is one regularized sequence per channel.  This step does
\emph{not} force different channels to share the same timestamps internally;
it only ensures that each channel is self-consistent on its own lattice.

\subsection{Anchor-Based Multi-Channel Sampling}
\label{appsub:anchor-based-multi-channel-sampling}

After regularization, a sample is built around a common anchor key $t$.  For
each channel separately, the dataloader selects the last record whose key is
less than or equal to $t$.  That selected record is the terminal record of the
input window on that channel.  The future window on the same channel begins at
the first record whose key is strictly greater than $t$.

Therefore the synchronization variable shared across channels is the anchor key
$t$ itself, not a row-by-row timestamp identity.  Two channels may have
different actual \field{close\_time} values at the same tensor position, as
long as those positions play the same role relative to the anchor key.

\subsection{Input and Future Windows}
\label{appsub:input-and-future-windows}

For each edge and channel, the input tensor $\tensor{x}_{u\,v}(t)$ is
back-aligned to the anchor:
it contains the past sequence ending at the last available record not after
$t$.  The future tensor $\tensor{f}(t)$ is forward-aligned from the same anchor:
it contains the sequence starting at the first available record after $t$.

At the concatenated-dataset level, anchor keys are restricted to the
intersection of the valid per-channel sampling domains.  This guarantees that
every retained anchor admits both an input window and a future window for every
active channel.

\subsection{Padding and Masks}
\label{appsub:padding-and-masks}

Different channels may contribute different native input and future lengths.
When channels are stacked into one sample:
\begin{enumerate}
\item Past windows are padded at the front, so the last valid input position
stays attached to the channel-specific anchor record.
\item Future windows are padded at the end, so the first valid future position
stays attached to the first channel record after the anchor.
\item The masks $\tensor{m}_x(t)$ and $\tensor{m}_f(t)$ mark which tensor positions
carry real records and which positions are only structural padding or invalid
nulls inherited from the lattice regularization step.
\end{enumerate}

This is why equal tensor indices across channels should be read as equal
relative positions with respect to the shared anchor key, not as equal raw
timestamps.
